% =====================================================================
%  PAPER IV -- from the CERTIFLIGHT concept (MSCA O1/O2, T1-T2, WP2/WP3).
%  Target: USENIX Security.
%
%  Standalone. Set as Overleaf's main document; pdfLaTeX + BibTeX.
%    pdflatex PaperIV_crosslayer_benchmark && bibtex PaperIV_crosslayer_benchmark \
%      && pdflatex PaperIV_crosslayer_benchmark && pdflatex PaperIV_crosslayer_benchmark
% =====================================================================
\documentclass[11pt,a4paper]{article}
% =====================================================================
%  Shared preamble for the three research proposals.
%
%  Each proposal \input{}s this file and is otherwise standalone, so any
%  one of them can be set as Overleaf's main document and compiled on its
%  own. Nothing here is specific to a single proposal.
% =====================================================================
\usepackage[T1]{fontenc}
\usepackage[utf8]{inputenc}
% Times, matching the parent manuscript. mathptmx is in every TeX Live and on
% Overleaf; lmodern is deliberately not used, since it is absent from minimal
% distributions and this preamble should compile anywhere.
\usepackage{mathptmx}
\usepackage[a4paper,margin=25mm]{geometry}
\usepackage{amsmath,amssymb,amsthm}
\usepackage{booktabs}
\usepackage{tabularx}
\usepackage{xcolor}
\usepackage{titlesec}
\usepackage{enumitem}
\usepackage{microtype}
\usepackage[hidelinks,breaklinks]{hyperref}

% ---- shared style tokens, matching the parent manuscript ------------
\definecolor{netcol}{RGB}{31,78,121}
\definecolor{autocol}{RGB}{18,120,90}
\definecolor{bridgecol}{RGB}{176,96,0}
\definecolor{accent}{RGB}{140,20,20}

\newcommand{\layernet}{\textcolor{netcol}{\textbf{network}}}
\newcommand{\layerauto}{\textcolor{autocol}{\textbf{autonomy}}}
\newcommand{\dataset}[1]{\textsc{#1}}
\newcommand{\code}[1]{\texttt{#1}}
\newcommand{\Real}{\mathbb{R}}
\newcommand{\Norm}[1]{\left\lVert #1\right\rVert}
\newcommand{\MCR}{\mathrm{MCR}}

\newtheorem{claim}{Claim}
\newtheorem{proposition}{Proposition}
\newtheorem{definition}{Definition}
\theoremstyle{remark}
\newtheorem{remark}{Remark}

% ---- section styling -------------------------------------------------
\titleformat{\section}{\large\bfseries\color{netcol}}{\thesection}{0.7em}{}
\titleformat{\subsection}{\normalsize\bfseries}{\thesubsection}{0.6em}{}
\titlespacing*{\section}{0pt}{1.4em}{0.6em}

\setlist[itemize]{leftmargin=1.35em,topsep=0.3em,itemsep=0.25em}
\setlist[enumerate]{leftmargin=1.6em,topsep=0.3em,itemsep=0.25em}

% ---- a boxed "what this rests on" callout ---------------------------
\newenvironment{honestly}
  {\par\medskip\noindent\begin{tabular}{@{}p{2.5mm}@{\hspace{3mm}}p{0.9\linewidth}@{}}
   \textcolor{accent}{\rule{2.5mm}{\baselineskip}} &
   \small\itshape}
  {\end{tabular}\par\medskip}

% ---- the proposal header --------------------------------------------
\newcommand{\proposalheader}[3]{%
  \begin{center}
  {\Large\bfseries #1\par}\medskip
  {\large\itshape #2\par}\medskip
  {\small #3\par}
  \end{center}
  \medskip\hrule\medskip}

\usepackage{array}
\newcolumntype{Y}{>{\raggedright\arraybackslash}X}
\newcolumntype{P}[1]{>{\raggedright\arraybackslash}p{#1}}
\newcommand{\dnet}{\delta_{\mathrm{net}}}
\newcommand{\dfault}{\delta_{\mathrm{fault}}}
\newcommand{\denv}{\delta_{\mathrm{env}}}
\newcommand{\dgnss}{\delta_{\mathrm{gnss}}}
\newcommand{\dtot}{\delta_{\mathrm{tot}}}
\newcommand{\prov}[1]{\textsc{\small[#1]}}
% Quantifier-level labels (defined in full in Paper II; referenced here).
\newcommand{\Qlvl}[1]{\textbf{Q}\textsubscript{\textit{#1}}}
\newcommand{\Qop}{\Qlvl{op}}
\newcommand{\Qpop}{\Qlvl{pop}}

\begin{document}

\proposalheader
  {Paper IV --- The Cross-Layer Evaluation}
  {When the Attacker Is Indistinguishable from the Weather: A Cross-Layer
   Benchmark and a Unified Perturbation Budget for Certified UAV Autonomy}
  {Research proposal from the CERTIFLIGHT concept (WP2/WP3) $\cdot$ RobustUAVs.ai
   $\cdot$ Roger Nick Anaedevha $\cdot$ 2026-10-01\\
   Target: USENIX Security $\cdot$ extends the parent manuscript~\cite{parentpaper}}

\begin{abstract}
\noindent
A UAV's navigation stack cannot tell, from the signal alone, whether its position
error is caused by a jammer, a drifting inertial unit, or a gust. We take that
indistinguishability as a security principle rather than a nuisance. If a
certificate is to mean anything to an operator, it must bound the mission outcome
under \emph{whatever} produced the perturbation, because the attacker is free to
choose the component that the defence models least well. We develop a unified
perturbation budget $\dtot=\dnet(\theta)\oplus\dgnss\oplus\dfault\oplus\denv$ that
aggregates an adversarial residual (a function of a network detector's operating
point), GNSS denial, stochastic faults and bounded environmental forcing into one
certified floor on the mission-completion rate --- and we show the security
consequence the single-source view misses: an adversary who shapes its
interference to \emph{mimic} a fault or a gust evades a detector tuned to
malicious signatures while producing the same navigation degradation. The
evaluation substrate is the missing ingredient in this field: a reproducible
cross-layer benchmark whose unit of record is the pairing of a network or
environmental event with the navigation window it induces, each pairing carrying
an auditable provenance flag that states whether the coupling was measured on one
platform, derived from a physical model, or aligned in simulation. We quantify
what the union-bound aggregation costs in tightness, and name the conditions
under which a joint treatment recovers it.
\end{abstract}

% =====================================================================
\section{The security argument}
% =====================================================================
Three communities each bound one layer and leave the interface to assumption.
Network intrusion detection bounds what reaches the aircraft but says nothing
about the flight after a missed intrusion. Certified
robustness~\cite{cohen2019smoothing,gronwall1919} bounds what a model does with a
perturbation but not where it comes from or how often an adversary can make one.
Fault-tolerant control rejects disturbances but has not been composed with
either. The property an operator, insurer and regulator actually need --- will
this mission complete, and can we prove it? --- falls between them.

The security content is sharper than ``let us unify three budgets''. It is an
\textbf{adversarial-equivalence} observation:

\begin{quote}\itshape
From the navigation input alone, an adversary's shaped interference, a sensor
fault, and an environmental disturbance are indistinguishable. An attacker will
therefore choose the disguise the defender models least well. A certificate that
is sound only against a declared attack model is unsound against the attacker who
dresses as the weather.
\end{quote}

\noindent
This is why the budget must be source-agnostic to be a security guarantee at all,
and why a detector tuned to malicious signatures is the wrong last line: the
decisive attack is the one that looks benign. A physical-invariant monitor such
as SAVIOR~\cite{quinonez2020savior} checks the dynamics rather than the signature
and is the right comparison point; our contribution is to carry its spirit
through to a \emph{certified mission floor} and to a benchmark that can measure
the evasion.

% =====================================================================
\section{The unified budget}
% =====================================================================
Navigation state $\dot x=F(x,u)$, $F$ Lipschitz with constant $L$ over the
operating region, horizon $T$. A bounded input perturbation $\|\delta u\|\le\dtot$
gives, through Gr\"onwall, a trajectory-tube radius $\rho=\dtot e^{LT}$; the
mission completes while the tube fits the safe-corridor margin $m$, so the
certified floor is $f(\dtot)=\Pr[\dtot e^{LT}\le m]$ (parent manuscript).

The new object is the budget itself:
\begin{equation}
\label{eq:budget}
\dtot \;=\; \dnet(\theta)\ \oplus\ \dgnss\ \oplus\ \dfault\ \oplus\ \denv ,
\end{equation}
aggregating sources with \emph{different probabilistic semantics}: $\dnet(\theta)$
deterministic worst-case (adversarial residual below detector threshold $\theta$,
bounded by the residual-budget lemma $\Delta(\theta)\le H\min(\theta,s)$),
$\dgnss$ and $\dfault$ stochastic, $\denv$ bounded-but-unknown.

\begin{proposition}[Heterogeneous composition, with its tightness cost]
\label{pr:compose}
The $\oplus$ of a deterministic, a stochastic and a bounded-unknown source admits
a joint certificate. The union bound over independently certified sources is
sound but loose; the looseness is exactly the dependence-uncertainty gap of
Makarov's bound~\cite{makarov1981,embrechts2013var} on the distribution of a sum
with fixed marginals and unknown coupling. A joint treatment recovers the union
bound's slack precisely when the sources are comonotone, and we characterise the
intermediate regime.
\end{proposition}

\begin{honestly}
Naming the tightness gap as the Makarov/VaR-aggregation gap is the theoretical
hinge, and it is a claim to prove, not assert: the mapping from ``three
perturbation sources with mixed semantics'' to ``sum of random variables with
fixed marginals and unknown copula'' has to be made rigorous, and the
deterministic source needs a degenerate-marginal treatment. If that reduction
does not hold cleanly, the paper falls back to the plain union bound --- which is
already sound --- and reports the gap empirically instead.
\end{honestly}

% =====================================================================
\section{The benchmark (the field's missing substrate)}
% =====================================================================
A guarantee spanning two layers needs data objects that record the coupling, and
no public corpus does. The benchmark's unit of record is the \textbf{cross-layer
pairing}: a network or environmental window joined to the navigation window it
induces, carrying the detector operating point, the residual budget, the induced
perturbation, and the mission outcome.

\paragraph{Provenance is mandatory and enforced.} Every pairing carries a flag:
\code{measured\_same\_platform} (both windows observed on one airframe),
\code{physics\_model} (coupling through a stated physical relation), or
\code{synthetic\_alignment} (aligned in simulation). A reviewer filters by
evidence strength rather than trusting an unaudited join. This is the discipline
that keeps the central honest qualification visible: of six corpora, only one
observes a network event and the flight it degrades on the same platform, and its
release carries no machine-readable attack intervals --- so the strongest
evidence class is sparse, and the benchmark says so in every record rather than
burying it.

\paragraph{Contents.} Lossless ingestion adapters for six public corpora; a
physics-informed navigation benchmark; a detector-operating-curve harness that
sweeps $\theta$ and reports recall and false-positive rate (the network-side
premise the composed guarantee consumes); and a reference pairing pipeline.
Third-party data is fetched from origin under its own licence; adapters and
unified labels are released, not redistributed corpora.

% =====================================================================
\section{The evasion study (the USENIX core)}
% =====================================================================
The benchmark makes the adversarial-equivalence claim measurable. The study:

\begin{description}[leftmargin=2.6em,style=nextline]
\item[RQ1] \textbf{Can an attacker mimic a fault/gust to evade a
  signature-tuned detector while producing the same $\delta$?} Construct shaped
  interference whose network-layer signature falls below $\theta$ but whose
  induced navigation degradation equals that of a modelled fault, and measure the
  detector's miss rate against the induced $\delta$.
\item[RQ2] \textbf{Does the composed certificate bound the mission outcome even
  when the detector is evaded?} This is the payoff: the floor should hold because
  the budget is source-agnostic, where a detector-only defence fails silently.
\item[RQ3] \textbf{How loose is the union bound in practice}, across the six
  corpora and the physics benchmark, and where does the joint treatment of
  Proposition~\ref{pr:compose} recover it?
\item[RQ4] \textbf{Soundness:} does the certified floor lower-bound the measured
  completion rate at \emph{every} operating point, across seeds, under
  Holm-corrected tests? A violation is a bug in the instantiation.
\end{description}

\paragraph{Why USENIX Security.} The contribution is a threat-model correction
(the attacker who dresses as the weather), an evaluation methodology with enforced
provenance, and a released artifact --- the shape of contribution USENIX rewards,
and the venue where an evaluation substrate can become a community default. The
theorem is the frame; the benchmark and the evasion measurement are the paper.

% =====================================================================
\section{What is in hand, what is new}
% =====================================================================
\begin{center}\small
\begin{tabularx}{\linewidth}{@{}P{0.40\linewidth}Y@{}}
\toprule
\textbf{In hand (parent + Papers I--II)} & \textbf{New here} \\
\midrule
Six-corpus ingest, certificate engine, four certificate methods &
$\dfault$ and $\denv$ as budget terms with mixed semantics \\
Residual-budget lemma; $\gamma=1.625$\,m/s campaign supremum &
the adversarial-equivalence threat model and the evasion study \\
Detector-operating-curve harness; provenance flags &
Proposition~\ref{pr:compose} and the Makarov/VaR tightness characterisation \\
\dataset{HCRL UAVCAN}, \dataset{Whelan} and four more adapters &
public benchmark release whose unit is the cross-layer pairing \\
\bottomrule
\end{tabularx}
\end{center}

% =====================================================================
\section{Relation to the other papers}
% =====================================================================
\begin{itemize}
\item \textbf{Parent manuscript} proves the composition for the time-delay class
  on the network side. Paper~IV broadens the \emph{source} of the perturbation to
  faults and environment and makes the threat model adversarial-equivalence; it
  is the security-and-measurement paper to the parent's theory paper.
\item \textbf{Paper~I} varies $\gamma$ and $\Delta$ as design variables; Paper~IV
  varies the \emph{source} of $\delta$. Complementary axes of the same budget.
\item \textbf{Paper~III} consumes the per-agent radius this budget produces.
  Paper~IV is upstream of it.
\item \textbf{Paper~II (SoK)} supplies the vocabulary; Paper~IV is a \Qop{}-level
  security evaluation within it.
\end{itemize}

% =====================================================================
\section{Risks}
% =====================================================================
\begin{itemize}
\item \textbf{The equivalence attack may be only partially realisable.} Shaping
  interference to match a fault's $\delta$ \emph{and} stay below $\theta$ may not
  be achievable for every fault profile. That is itself a finding --- the subset
  of degradations an attacker can disguise is the real attack surface --- and the
  paper reports the achievable region rather than claiming the whole.
\item \textbf{The tightness reduction (Prop.~\ref{pr:compose}) may not close.}
  Fallback: the sound union bound plus an empirical gap measurement. The benchmark
  and evasion study stand regardless.
\item \textbf{Same-platform pairings stay scarce.} The provenance flags make this
  explicit rather than hidden; \code{physics\_model} and
  \code{synthetic\_alignment} pairings sustain coverage, and the limitation is
  reported as a limitation. This is a standing constraint from the parent work,
  carried forward honestly.
\item \textbf{Fault/gust realism.} $\dfault$ and $\denv$ are modelled from the
  fault-tolerant-control literature; every modelled number carries its provenance
  tag, and the paper does not present a simulated gust as a measured one.
\end{itemize}

% =====================================================================
\section{Page plan and timeline}
% =====================================================================
\begin{center}\small
\begin{tabular}{@{}lr@{\qquad}lr@{}}
\toprule
\textbf{Section} & \textbf{pp} & \textbf{Section} & \textbf{pp} \\
\midrule
Introduction                 & 1.0  & Evasion study and results     & 2.5 \\
Threat model (equivalence)   & 1.25 & Soundness and union-bound gap & 1.5 \\
Unified budget and theory    & 1.75 & Discussion, ethics, limits    & 1.0 \\
Benchmark design             & 1.75 & Related work, conclusion      & 1.0 \\
\midrule
\multicolumn{3}{@{}l}{\textbf{Body (USENIX 13-page limit)}} & \textbf{11.75} \\
\bottomrule
\end{tabular}
\end{center}

\noindent
About five to six months. The ingest, harness and certificate engine exist; the
new work is the fault and environment budget terms, the evasion construction, and the
benchmark packaging and release. An ethics subsection is mandatory: the evasion
method is dual-use, and following the parent's evasion analysis, it is presented
with its mitigation (invariant-based monitoring, and source-agnostic
certification) in the same section, bounds what the attacker achieves, and
targets a defensive artifact.

\small
\bibliographystyle{plain}
\bibliography{proposals}

\end{document}
